\chapter{证据驱动的声明验证方法}\label{ch:openbook}

第~\ref{ch:clara}~章表明，在不引入外部信息的条件下，把 LLM 的分析做得更细并不能改变检测器主要依赖文风这一事实。本章转向\emph{开卷}设定：每条待核查的声明附带一组从网络检索到的原始报道，模型需要从中找出证据并据此判断声明的真实性等级。本章要回答三个问题：证据究竟能带来多大的提升；现有数据集中的“证据”是否泄漏了结论；一个可在单张显卡上训练的 LLM 验证器，能否既用好证据，又在证据不可用时不被误导。

本章的安排如下：\ref{sec:ob-task}~节给出任务定义与数据集；\ref{sec:ob-audit}~节审计证据来源并定义评测协议；\ref{sec:ob-method}~节提出证据自筛选验证方法 \ours{}；\ref{sec:ob-setup}~节说明实验设置；\ref{sec:ob-results}~节报告并分析实验结果；\ref{sec:ob-discussion}~节讨论局限；\ref{sec:ob-summary}~节小结。

\section{任务定义与数据集}\label{sec:ob-task}

\subsection{任务定义}

给定一条声明 $c$ 以及为它检索到的原始报道集合 $\mathcal D_c=\{d_1,\dots,d_{n_c}\}$，每篇报道 $d$ 是一个句子序列并带有来源网址。任务是预测 $c$ 的真实性等级 $y\in\{1,\dots,C\}$，等级按“从最不属实到最属实”排序。除等级外，我们还希望模型（i）指出它依据的证据句，（ii）在检索到的材料与声明无关时能够发现这一点。

训练数据除标签外还带有两类\emph{特权信息}：事实核查机构为该声明撰写的结论文本 $r_c$，以及数据集构建者依据 $r_c$ 为报道中每个句子 $s$ 标注的\emph{银标证据标签} $z(s)\in\{0,1\}$。两者都来自事实核查文章，在声明尚未被核查的真实场景中并不存在，因此本章只把它们用作训练信号，任何测试时的输入都不包含它们。

\subsection{数据集}

实验使用 Yang 等\citep{yang2022cofced}构建的 RAWFC 与 LIAR-RAW。RAWFC 的声明与结论来自 Snopes，分为假、半真、真三类；LIAR-RAW 在 LIAR\citep{wang2017liar}的基础上为每条声明补充了原始报道，声明与结论来自 PolitiFact，沿用其六级标签（pants-fire、false、barely-true、half-true、mostly-true、true）。两个数据集都使用原作者发布的训练、验证与测试划分。按原论文的说明，报道是用声明的关键词经搜索引擎检索得到的（每条声明至多 30 篇），银标证据标签由句子与核查结论文本之间的 ROUGE 与语义相似度超过给定阈值得到\citep{yang2022cofced}。表~\ref{tab:ob-data} 给出统计信息。

\begin{table}[htbp]
\centering\small
\caption{RAWFC 与 LIAR-RAW 的统计信息}\label{tab:ob-data}
\fittable{gen/ob_tables/data.tex}
\note{候选句为一条声明的全部报道中去重后不少于 5 个词的句子；“严格协议”见 \ref{sec:ob-audit}~节。}
\end{table}

两个数据集的差别值得注意。RAWFC 的规模小（测试集只有 \cnt{R/stats/test/claims} 条），但每条声明的候选句多，其中约四分之一带有银标证据标签；LIAR-RAW 的测试集有 \cnt{L/stats/test/claims} 条，声明多为政治人物的简短表态，候选句中银标证据不到十分之一，并且需要区分六个相邻的等级。前者的结果受随机波动影响更大，后者的任务本身更难。

\section{证据来源的审计与评测协议}\label{sec:ob-audit}

\subsection{证据泄漏问题}

开卷评测隐含一个前提：证据应当是在声明被核查\emph{之前}就能得到的材料。如果检索到的报道本身就是一篇事实核查文章，或者转述了核查结论，那么“验证”就退化为抄写结论，评测结果不再反映模型的核查能力\citep{glockner2022missing}。RAWFC 与 LIAR-RAW 的构建者意识到这一点，在收集报道时过滤了事实核查网站，并去掉了发布时间晚于核查文章的报道\citep{yang2022cofced}。此后在这两个数据集上的工作大多直接使用全部报道；FFRR\citep{zhang2024ffrr}用 Glockner 等\citep{glockner2022missing}的检测方法进一步去除了被判为泄漏的文档（RAWFC 中占 4.7\%，LIAR-RAW 中占 2.4\%），但没有给出去除前后的对照。

\subsection{审计方法与结果}

我们对两个数据集的全部报道做了一次基于规则的审计，按以下五种原因标记可能泄漏结论的报道：（1）来自标签来源的核查机构（Snopes 或 PolitiFact）；（2）来自其他专门的事实核查网站，或事实核查文章的聚合页面；（3）来自代码与数据托管站点（如 GitHub、Kaggle），这类页面往往是基准数据集本身的副本；（4）来自一般新闻媒体，但网址表明它属于事实核查栏目；（5）不属于以上各类，但正文提及事实核查（例如出现“fact check”“Snopes”“PolitiFact”“debunk”等字样）。前四种按网址判断，第五种按正文判断。完整的规则见附录~\ref{app:ob-impl}。

\begin{table}[htbp]
\centering\small
\caption{两个数据集中被标记的报道及其涉及的声明}\label{tab:ob-audit}
\fittable{gen/ob_tables/audit.tex}
\note{“报道数”“涉及声明”统计全部三个划分，“测试集声明”只统计测试集。一篇报道可以同时因多种原因被标记；“合计”为至少因一种原因被标记。}
\end{table}

表~\ref{tab:ob-audit} 给出审计结果，有三点值得说明。

第一，被标记的报道占比不高，但涉及的声明不少。RAWFC 中被标记的报道占 \pone{R/audit/total/reports_any_share}\%，LIAR-RAW 中占 \pone{L/audit/total/reports_any_share}\%；由于每条声明有十几到二十几篇报道，至少有一篇报道被标记的声明分别占 \pone{R/audit/total/claims_any_share}\% 与 \pone{L/audit/total/claims_any_share}\%，在 LIAR-RAW 的测试集中达到 \pone{L/audit/test/claims_any_share}\%。

第二，LIAR-RAW 的报道中混有\emph{数据集本身的副本}。有 \cnt{L/audit/total/reports_by_flag/dump} 篇报道来自代码与数据托管站点，涉及 \cnt{L/audit/total/claims_by_flag/dump} 条声明，其中多数是他人上传到 GitHub 的 LIAR 数据文件。这类页面逐行列出声明及其标签。在测试集中，有 \cnt{L/dump_rows/test/claims_with_dump} 条声明带有这类报道；数据集只保留了每个页面的少量句子，其中 \cnt{L/dump_rows/test/own_row} 条声明自己所在的那一行恰好被保留下来，且六级标签就印在声明文本之前。能被直接读出标签的声明很少，但这一现象说明，按域名过滤核查网站并不能挡住所有泄漏途径。在我们查阅的范围内，尚未见到已有工作指出这一点。

第三，这一审计是保守的启发式规则，既有误报也有漏报。第五种原因会标记一些只是顺带提到“事实核查”的普通报道；而转述了核查结论、却不含任何关键词的报道不会被标记\citep{glockner2022missing}。数据集没有提供报道的发布时间，我们也无法检查它是否晚于核查文章。因此，下面定义的严格协议应理解为“去掉了能够识别的泄漏”，而不是“保证没有泄漏”。

\subsection{两种评测协议}

据此我们定义两种证据协议。\emph{全部报道}协议使用数据集提供的全部报道，这是已有工作的通行做法。\emph{严格}协议去掉被标记的全部报道，只从剩余报道中选取证据。本章以严格协议为默认设置（训练与测试都只使用未被标记的报道），同时报告同一模型在全部报道协议下的测试结果，并在 \ref{sec:ob-leak}~节专门比较两者。严格协议下，LIAR-RAW 的测试集中有 \cnt{L/stats/test/pool_strict/empty} 条声明没有任何候选句，模型此时被告知“没有检索到证据”。

\section{证据自筛选验证方法 \ours{}}\label{sec:ob-method}

\subsection{设计思路}

第~\ref{ch:clara}~章的结果对开卷方法的设计有三点直接的启示。其一，直接微调的 LLM 是比“冻结的 LLM 加小模型”更强的骨干，因此本章以 LoRA 微调的 LLM 作为验证器，而不是让小模型去阅读 LLM 生成的分析。其二，用错配负例训练出来的“这段分析是否属于这条新闻”的判断是可靠的；检索到的证据同样可能与声明对不上，这一做法可以迁移过来。其三，方法声称利用了什么信息，就应当通过改变这一信息的测试来验证。

直接把检索到的句子拼入提示并微调，是最简单的开卷验证器。它有两个不足。第一，它能读入的句子数受显存限制，读入哪些句子完全由第一阶段的检索器决定，而词项匹配检索只衡量句子与声明在用词上的重合，不衡量句子对判断真伪是否有用。第二，训练时它见到的总是属于该声明的证据，从未被要求判断“证据是否可用”；\ref{sec:ob-robust}~节将显示，当检索返回的是无关材料时，它会以很高的置信度把几乎所有声明判为假。

\ours{}（\textbf{S}elect\textbf{I}ve \textbf{EV}idence v\textbf{E}rifier）针对这两点做了三处设计（图~\ref{fig:ob-framework}）：让验证器在判断真伪的同一次前向计算中为每个证据句打分；训练时混入错配的证据，并在每轮重新抽取证据；推理时先读、再选、后判。整个方法只包含一个 LoRA 微调的 LLM 与一个线性的打分头，推理时不生成任何文本。

\begin{figure}[htbp]
\centering
\resizebox{\linewidth}{!}{% Framework figure of the open-book method (TikZ); natural width ~180 mm, scaled to \linewidth.
% Blue = the verifier (one LoRA-tuned LLM, used in both passes), green = selection,
% orange = what exists only during training.
\definecolor{cLLM}{HTML}{EB6834}
\definecolor{cSLM}{HTML}{2A78D6}
\definecolor{cGate}{HTML}{1BAF7A}
\definecolor{cInk}{HTML}{0B0B0B}
\definecolor{cMuted}{HTML}{6B6A66}
\definecolor{cLine}{HTML}{C3C2B7}
\begin{tikzpicture}[
  x=1mm, y=1mm, font=\footnotesize, >={Stealth[length=1.6mm]},
  box/.style={draw=cLine, rounded corners=2pt, fill=white, align=center, inner sep=2.5pt, text=cInk},
  ver/.style={box, draw=cSLM!80, fill=cSLM!7},
  sel/.style={box, draw=cGate!90, fill=cGate!9},
  trn/.style={box, draw=cLLM!80, fill=cLLM!6, align=left},
  flow/.style={->, draw=cMuted, line width=0.5pt},
  sup/.style={->, draw=cMuted, dashed, line width=0.5pt},
  lab/.style={font=\scriptsize, text=cMuted, inner sep=1.5pt}]

% ---------------- inference: retrieve, read, select, verify
\node[box, minimum width=24mm] (c) at (12,0) {声明 $c$\\[-1pt]{\scriptsize 原始报道 $\mathcal D_c$}};
\node[box, minimum width=30mm] (pool) at (46,0) {候选池 $\mathcal P_c$\\[-1pt]{\scriptsize 协议过滤、去重}\\[-1pt]{\scriptsize BM25 取前 $K$ 句}};
\node[ver, minimum width=40mm] (read) at (88,0) {\textbf{读}：验证器 $f_\theta$\\[-1pt]{\scriptsize 每次读 $m$ 句，共 $\lceil K/m\rceil$ 次}\\[-1pt]{\scriptsize 句末隐状态 $\to$ 证据得分 $u(s)$}};
\node[sel, minimum width=28mm] (sel) at (128,0) {\textbf{选}\\[-1pt]{\scriptsize 得分最高的}\\[-1pt]{\scriptsize $m$ 句 $E^{*}$}};
\node[ver, minimum width=40mm] (ver) at (168,0) {\textbf{判}：验证器 $f_\theta$\\[-1pt]{\scriptsize 读入 $c$ 与 $E^{*}$}\\[-1pt]{\scriptsize 选项字母的 logit $\to p_\theta(y\mid c,E^{*})$}};
\draw[flow] (c) -- (pool);
\draw[flow] (pool) -- (read);
\draw[flow] (read) -- (sel);
\draw[flow] (sel) -- (ver);

% ---------------- outputs
\node[box, minimum width=34mm] (use) at (128,-23) {可用性得分 $a(c)$\\[-1pt]{\scriptsize $E^{*}$ 上证据得分的均值}};
\node[box, minimum width=22mm] (out) at (176,-23) {真实性等级 $\hat y$};
\draw[flow] (sel) -- (use);
\draw[flow] ([xshift=8mm]ver.south) -- (out.north);

% ---------------- training only
\node[trn, minimum width=62mm] (train) at (33,-47) {%
  {\scriptsize \textbf{训练查询的证据 $E$}（每轮重新构造）}\\[-1pt]
  {\scriptsize 重采样：从 $\mathcal P_c$ 按名次加权抽取 $m$ 句}\\[-1pt]
  {\scriptsize 错配：以概率 $\rho$ 换成另一条声明的候选句}};
\node[trn, minimum width=48mm] (lev) at (93,-47) {%
  {\scriptsize $\lambda_{\mathrm{ev}}\,\overline{\mathrm{BCE}}\big(u_i,\,z(e_i)\big)$}\\[-1pt]
  {\scriptsize $z$：银标证据标签}\\[-1pt]
  {\scriptsize 错配时全为 0}};
\node[trn, minimum width=56mm] (lce) at (152,-47) {%
  {\scriptsize $\mathrm{CE}\big(p_\theta(\cdot\mid c,E),\,y\big)$}\\[-1pt]
  {\scriptsize $y$：真实性等级标签}\\[-1pt]
  {\scriptsize 错配时仍为本声明的标签}};
\draw[sup] ([xshift=-5mm]lev.north) -- ([xshift=0mm]read.south);
\draw[sup] ([xshift=0mm]lce.north) -- ([xshift=-16mm]ver.south);
\begin{scope}[on background layer]
  \node[draw=cLine, dashed, rounded corners=3pt, inner sep=3pt, fit=(train)(lev)(lce)] (plate) {};
\end{scope}
\node[lab, anchor=south west] at ([xshift=1mm,yshift=0.3mm]plate.north west) {仅在训练时};
\end{tikzpicture}
}
\caption{\ours{} 的总体框架}\label{fig:ob-framework}
\note{上半部分为推理流程：从候选池读入证据句并打分，选出得分最高的句子，再据此判断。“读”与“判”使用同一个验证器。下半部分为只在训练时使用的证据构造方式与损失。}
\end{figure}

\subsection{候选池与第一阶段检索}

对声明 $c$，把评测协议允许使用的报道拆成句子，去掉重复的句子与少于 5 个词的片段，得到候选句集合。数据集发布的报道是经过词形还原的文本，而声明不是，因此检索前对两者都做词干提取。我们用 BM25（式~\eqref{eq:bm25}）以声明为查询给候选句打分，文档频率在该声明自己的候选句集合内统计，取得分最高的 $K$ 句构成\emph{候选池} $\mathcal P_c=(s_{(1)},\dots,s_{(K)})$。第一阶段只负责缩小范围，不需要训练。

\subsection{验证器}

验证器 $f_\theta$ 是一个指令微调的 LLM。它的输入由系统提示、声明、带编号的证据句列表以及一个选择题组成：证据句按报道分组（同一报道的句子保持原有顺序），每句标明来源域名；选择题列出全部 $C$ 个真实性等级，每个等级对应一个字母，并要求只用一个字母作答。提示词全文见附录~\ref{app:ob-prompt}。类别分布按 \ref{sec:bg-readout}~节的方式读出：在回答的起始位置只取 $C$ 个选项字母的 logit 并归一化，记为 $p_\theta(y\mid c,E)$，其中 $E=(e_1,\dots,e_m)$ 是读入的证据句。验证器用 LoRA 微调，损失为该分布与标签的交叉熵。

不给证据列表、只给声明时，同一个验证器就是闭卷的 LoRA-LLM，即第~\ref{ch:clara}~章中最强的基线；不做微调时，它是零样本的 LLM。这三者在本章都作为对比方法出现，它们与 \ours{} 使用完全相同的读出方式。

\subsection{证据打分头}\label{sec:ob-head}

记 $t_i$ 为证据句 $e_i$ 的最后一个词元在输入中的位置，$\bm h_{t_i}$ 为最后一层在该位置的隐状态。我们加一个线性的打分头，
\begin{equation}
u_i=\sigma\big(\bm w^{\top}\bm h_{t_i}+b\big)\in(0,1),\label{eq:ob-head}
\end{equation}
把 $u_i$ 称为证据句 $e_i$ 的\emph{证据得分}。由于声明排在证据之前，而解码器的注意力是单向的，$\bm h_{t_i}$ 已经“读过”声明与这句证据，却看不到排在后面的内容。因此，一次前向计算同时给出全部 $m$ 个证据得分与类别分布，打分不需要额外的模型，也不需要额外的调用。

证据得分用银标证据标签监督：
\begin{equation}
\mathcal L_{\mathrm{ev}}=\frac{1}{m}\sum_{i=1}^{m}\operatorname{BCE}\big(u_i,\,z(e_i)\big).
\end{equation}
银标标签来自句子与核查结论文本的相似度，它告诉模型的是“核查人员最终用到了哪些内容”。这是一种只在训练时可用的特权信息；测试时模型只能依靠自己学到的打分。

\subsubsection{为什么可能有效，以及可能的失败方式}
让同一个模型既打分又判断，使“哪些句子有用”的监督直接作用在验证器的表示上；与单独训练一个重排序模型相比，它不增加参数与调用。潜在的问题有三个。银标标签是由相似度阈值自动得到的，含有噪声，模型可能学到“像核查文章的句子”而不是“对判断有用的句子”。句子的得分依赖于排在它前面的句子，同一句话在不同的上下文中得分可能不同。LIAR-RAW 中银标证据很稀疏，监督信号较弱。这些问题的实际影响在 \ref{sec:ob-evidence}~节检验。

\subsection{训练时的证据构造}\label{sec:ob-train}

\subsubsection{证据重采样}
如果训练时总是读入 BM25 排在最前面的 $m$ 句，打分头就只见过这些句子，而推理时它需要给整个候选池打分。因此在每一轮训练中，我们为每条声明从候选池中不放回地抽取 $m$ 句作为证据，名次为 $j$（从 0 起计）的句子被抽中的概率正比于 $1/(j+8)$：排名靠前的句子更常被选中，但池中每一句都有机会出现。重采样同时起到数据增广的作用，同一条声明在不同轮次搭配不同的证据。它的代价是，训练时读到的证据平均而言比推理时选出的证据更差。

\subsubsection{错配证据}
在每一轮训练中，每条声明以概率 $\rho$ 被换上\emph{另一条}训练声明的证据（由一个没有不动点的随机置换指定）。对这样的样本，真实性等级的目标仍然是该声明自己的标签，所有证据句的证据得分目标为 $0$。前一个目标要求验证器在证据与声明无关时不去“解读”这些证据，而是退回到只凭声明本身作出的判断；后一个目标让打分头学会识别不属于该声明的句子。

这与第~\ref{ch:clara}~章中相关性门控的训练方式一脉相承。它的局限也是相同的：随机抽取的另一条声明在话题上通常与原声明无关，模型学到的首先是话题层面的不匹配。对于话题相近、但并不针对该声明的材料，它能否识别需要专门检验，我们在 \ref{sec:ob-robust}~节用“相近声明的证据”这一条件考察。另外，退回到的闭卷判断本身主要依赖文风（第~\ref{ch:clara}~章），它只是“没有可用证据时能做到的最好结果”，并不可靠。

\subsubsection{训练目标}
总损失为
\begin{equation}
\mathcal L=\operatorname{CE}\big(p_\theta(\cdot\mid c,E),\,y\big)+\lambda_{\mathrm{ev}}\,\mathcal L_{\mathrm{ev}},
\end{equation}
其中 $E$ 是按上述方式为该样本构造的证据。可训练的参数只有 LoRA 的低秩矩阵与打分头。

\subsection{读—选—判推理}\label{sec:ob-infer}

推理分三步。

\textbf{读}：把候选池的 $K$ 个句子分成 $\lceil K/m\rceil$ 组，每组不超过 $m$ 句。分组时按名次轮流分配（名次为 $j$ 的句子进入第 $j \bmod \lceil K/m\rceil$ 组），使每组都同时含有排名靠前与靠后的句子，与训练时重采样得到的证据在构成上相近。验证器对每组做一次前向计算，得到池中每个句子的证据得分 $u(s)$。

\textbf{选}：取证据得分最高的 $m$ 句，记为 $E^{*}$。

\textbf{判}：验证器读入声明与 $E^{*}$，输出 $p_\theta(y\mid c,E^{*})$，取概率最大的等级为预测 $\hat y$。

此外，我们把最高的 $m$ 个证据得分的均值
\begin{equation}
a(c)=\frac{1}{m}\sum_{s\in E^{*}}u(s)
\end{equation}
作为声明级的\emph{证据可用性得分}：它很低时，说明候选池中没有模型认为有用的句子。

整个推理需要 $\lceil K/m\rceil+1$ 次前向计算，不生成任何文本。与只读 BM25 前 $m$ 句的单遍推理相比，它让验证器读过了整个候选池，并且最终依据的是它自己挑出的证据；代价是数倍的推理开销，以及打分头的错误会传递到最终判断。算法~\ref{alg:ob} 给出完整流程。

\begin{algorithm}[htbp]
\caption{\ours{} 的训练与推理}\label{alg:ob}
\KwIn{训练集 $\{(c,\mathcal D_c,y,z)\}$；候选池大小 $K$；每次读入的句数 $m$；错配概率 $\rho$；权重 $\lambda_{\mathrm{ev}}$}
\KwOut{验证器参数 $\theta$（LoRA 与打分头）}
对每条声明用 BM25 构造候选池 $\mathcal P_c$\;
\For{每一轮}{
  \ForEach{训练声明 $c$}{
    以概率 $\rho$ 取另一条声明 $c'$ 的候选池作为来源（错配），否则取 $\mathcal P_c$\;
    从来源中按名次加权抽取 $m$ 句得到 $E$；错配时令所有 $z(e_i)=0$\;
    一次前向计算得到 $p_\theta(\cdot\mid c,E)$ 与 $u_1,\dots,u_m$\;
    用 $\operatorname{CE}(p_\theta,y)+\lambda_{\mathrm{ev}}\mathcal L_{\mathrm{ev}}$ 更新 $\theta$\;
  }
  在验证集上用读—选—判推理计算 macro-F1，保留最优参数\;
}
\BlankLine
\textbf{推理}（声明 $c$）：分 $\lceil K/m\rceil$ 次读入 $\mathcal P_c$，得到每句的 $u(s)$\;
$E^{*}\leftarrow$ 得分最高的 $m$ 句；$\hat y\leftarrow\arg\max_y p_\theta(y\mid c,E^{*})$；$a(c)\leftarrow E^{*}$ 上 $u$ 的均值\;
\end{algorithm}

\subsection{与相关工作的区别}

让 LLM 自己参与对检索内容的判断并非新想法。Self-RAG\citep{asai2024selfrag}训练模型生成反思词元来评价检索到的段落，RankRAG\citep{yu2024rankrag}用同一个指令微调的 LLM 完成上下文排序与答案生成。\ours{} 的打分不经过生成：证据得分直接从每个证据句末尾的隐状态读出，一次前向计算即可为一组句子打分，并且与真实性判断共用这次计算。在编码器模型上，联合学习证据选择与真实性判断已有先例\citep{fajcik2022claimdissector,wadden2021multivers}，CofCED\citep{yang2022cofced}也用银标证据标签监督它的证据蒸馏模块；\ours{} 把这种联合监督放进了一个生成式 LLM 验证器，并用它驱动推理时的证据选择。

在训练中改变证据以约束模型对证据的依赖，与 REAL\citep{deng2026real}的思路相近。REAL 删去标注的证据句，并要求模型在证据被删去后改判为“证据不足”。本章使用的数据集没有“证据不足”这一类别，也没有可靠的证据标注；\ours{} 换上的是另一条声明的证据，要求模型此时退回到只凭声明的判断，并通过证据得分表明证据不可用。Atanasova 等\citep{atanasova2022insufficient}在编码器模型上研究了证据不足的识别。

与在同一数据集上的方法相比，L-Defense\citep{wang2024ldefense}及其后续工作\citep{wang2026gdefense}让 LLM 生成解释，再由小模型判断；DeReC\citep{qazi2025derec}用稠密检索加编码器分类器；直接微调 LLM 阅读证据的工作\citep{cheung2023factllama,kumar2025entailed,wang2025evicheck}没有联合的证据打分，也没有考察证据不可用的情形。

\section{实验设置}\label{sec:ob-setup}

\subsection{对比方法}

为了把“证据”“微调”与“本章的设计”三者的作用分开，我们自己实现了下列方法，它们使用同一个 LLM（Qwen2.5-7B-Instruct）、同一个第一阶段检索器与同一种证据协议：
\begin{itemize}
\item \textbf{零样本 LLM}：不做任何训练，分别在只给声明与给出 BM25 前 $m$ 句证据两种输入下读出类别分布。
\item \textbf{RoBERTa-large}\citep{liu2019roberta}：全量微调的编码器分类器，输入为声明，或声明与 BM25 前 $m$ 句证据的拼接（超过 512 个词元的部分截断）。它代表“检索加编码器分类器”这一类方法\citep{qazi2025derec}。
\item \textbf{LoRA-LLM}：LoRA 微调的验证器，分别在只给声明（闭卷）与给出 BM25 前 $m$ 句证据两种输入下训练与测试。后者是最直接的开卷验证器，也是 \ours{} 去掉全部三处设计后的形式。
\item \textbf{\ours{}}：本章提出的方法。
\end{itemize}
此外，我们列出已有工作在这两个数据集上报告的数值作为参照。这些数值取自各自的论文，是在全部报道协议下得到的，多数来自单次运行，所用的 LLM 也各不相同，与本章的结果不能做严格的比较。

\subsection{实现与超参数}

每次读入的证据句数 $m=16$，候选池大小 $K=48$，因而读—选—判推理需要 4 次前向计算。输入的最大长度为 1{,}100 个词元，超出时从证据列表的末尾去掉句子，声明与选择题不会被截断。LoRA 的秩为 16，作用于注意力与前馈层的全部线性投影；学习率 $10^{-4}$（打分头为 $10^{-3}$），线性预热后线性衰减，等效批大小 16。RAWFC 训练 4 轮，LIAR-RAW 的训练集约为其 6 倍，训练 2 轮；每轮结束后在验证集上评估，保留验证集 macro-F1 最高的参数，测试集只在选定的参数上评估一次。训练使用梯度检查点以容纳带证据的长输入。\ours{} 的两个超参数取预先设定的值，$\lambda_{\mathrm{ev}}=0.5$，$\rho=0.15$，没有做搜索。RoBERTa-large 在学习率 $\{1\times10^{-5},2\times10^{-5}\}$ 上各运行一遍，按种子平均的验证集 macro-F1 选择，RAWFC 训练 10 轮，LIAR-RAW 训练 5 轮。全部实验在单张 RTX 3090 上完成，其余细节见附录~\ref{app:ob-impl}。

\subsection{评测协议}

主指标为 macro-F1，同时报告宏平均精确率、召回率与准确率。可训练方法在 RAWFC 上运行 5 个随机种子，在 LIAR-RAW 上运行 3 个；消融与对照实验为 3 个种子。显著性判断沿用 \ref{sec:bg-task}~节的标准：对测试样本的配对 bootstrap 与对随机种子的 Welch 检验都在 0.05 水平上成立，才称差异显著。RAWFC 的测试集只有 \cnt{R/stats/test/claims} 条，其上几个点以内的差异往往达不到这一标准，解读时以 LIAR-RAW 上的结论为主。

需要说明两点。其一，已有工作报告的 F1 有两种口径：多数是各类别 F1 的平均，即本文的 macro-F1；另有一部分是由宏平均精确率与召回率求调和平均。两者在精确率与召回率相差较大时并不相同，本章在引用已发表数值时标明了口径。其二，在确定最终的实验配置之前，我们在 RAWFC 上用单个随机种子做过一次先导实验，并查看了它的测试集结果；\ours{} 的超参数在先导实验之前设定，之后没有调整，本章报告的全部结果来自此后重新运行的实验。

\subsection{证据受控的测试条件}\label{sec:ob-conditions}

除标准测试外，我们在测试时改变候选池，考察各方法的行为。标签始终不变，每种方法沿用它自己的推理方式。
\begin{itemize}
\item \textbf{缺少前排}：从候选池中去掉 BM25 排在最前面的 $m$ 句，模拟检索质量较差的情形。
\item \textbf{混入一半}：候选池的一半是该声明自己的句子，另一半来自随机抽取的另一条测试声明，两者按名次交错排列。
\item \textbf{相近声明}：候选池全部来自与该声明\emph{文本最相似}的另一条测试声明，模拟检索到了话题相关、但并不针对该声明的材料。
\item \textbf{随机声明}：候选池全部来自随机抽取的另一条测试声明，模拟检索完全失败。
\item \textbf{无证据}：不给任何证据，只给声明。
\end{itemize}
在后三种条件下，证据与声明无关，一个合理的验证器应当退回到与闭卷模型相当的水平，而不是被无关材料带偏。LIAR-RAW 上这些条件在测试集的一个固定的随机子集（400 条）上评估，与之比较的原始结果取同一子集。

另外设两项对照。\emph{只读证据}的对照在训练与测试时都隐去声明文本，只给证据：如果它与完整模型一样好，说明证据本身就泄漏了结论，或者模型并没有在做声明与证据之间的比对\citep{hansen2021reason}。\emph{证据上限}在测试时用银标证据标签从候选池中挑选证据（先取带标签的句子），它使用了测试样本的特权信息，不是合法的测试条件，只用来估计证据选择还有多大的改进空间。

\section{实验结果与分析}\label{sec:ob-results}

（待实验完成后撰写）

\section{讨论与局限}\label{sec:ob-discussion}

（待实验完成后撰写）

\section{本章小结}\label{sec:ob-summary}

（待实验完成后撰写）
